\section{Differential Path Contribution Maintenance}
\label{sec:dpcm}

The \cpi introduced in the previous section resolves the \emph{representation} challenge by encoding the $s$-clique space without explicit enumeration.
%
The remaining challenge is algorithmic: after a batch of minimum-support $r$-cliques is peeled, how should the supports of the surviving $r$-cliques be updated \emph{exactly} without reconstructing the affected local clique space from scratch?
%
In this section, we expose a higher-level view of our method.
%
We show that exact peeling on \cpi can be understood as \emph{differential path contribution maintenance}: each path contributes a local amount of support, and peeling only changes the contributions of paths incident to removed $r$-cliques.

\begin{definition}[Active path state]\label{def:path-state}
Consider the peeling process at round $t$ and let $\cpi_t$ denote the current \cpi.
For each active path $\TPath \in \cpi_t$, let $\sigma_t(\TPath)$ denote the local state maintained for $\TPath$.
%
This state records exactly the information needed to determine which active $s$-cliques are still encoded by $\TPath$ and how they contribute to the active $r$-cliques covered by $\TPath$.
\end{definition}

\begin{definition}[Path contribution]\label{def:path-contribution}
Let $\mathcal{K}_r^{(t)}$ and $\mathcal{K}_s^{(t)}$ be the active $r$-cliques and $s$-cliques at round $t$.
For an active $r$-clique $R \in \mathcal{K}_r^{(t)}$ and an active path $\TPath \in \cpi_t$, define the \emph{path contribution} of $\TPath$ to $R$ as
\[
\phi_t(R,\TPath)
=
\bigl|\{\,S \in \mathcal{K}_s^{(t)} \mid R \subseteq S,\; S \encodedby \TPath \,\}\bigr|.
\]
\end{definition}

Intuitively, $\phi_t(R,\TPath)$ counts how many currently valid $s$-cliques encoded by $\TPath$ contain $R$.
%
The support of $R$ is then obtained by summing these local contributions over all active paths.

\begin{theorem}[Support decomposition by path contributions]\label{thm:dpcm-support}
For every round $t$ and every active $r$-clique $R \in \mathcal{K}_r^{(t)}$,
\[
\deg^{G_t}_{(s)}(R)
=
\sum_{\TPath \in \cpi_t}\phi_t(R,\TPath).
\]
\end{theorem}

\begin{proof}
By Lemma~\ref{lem:clique_count}, every active path encodes a well-defined set of active $s$-cliques, and each such $s$-clique contributes one unit of support to every contained $r$-clique.
%
Moreover, the \cpi path semantics ensure that each active $s$-clique is represented by exactly one path.
%
Therefore, counting active $s$-cliques containing $R$ is equivalent to summing the numbers contributed by all active paths.
\end{proof}

\stitle{Differential update view.}
Suppose that at round $t$ the algorithm removes a batch $B_t \subseteq \mathcal{K}_r^{(t)}$ of minimum-support $r$-cliques.
%
Only paths that contain removed cliques, or are created by splitting such paths, can change their contribution values.
%
This yields a purely local update rule.

\begin{definition}[Affected paths]\label{def:affected-paths}
Let $\Gamma_t(B_t)$ denote the set of paths that either:
(i) contain some removed $r$-clique in $B_t$, or
(ii) are newly created when such paths are updated or split.
\end{definition}

\begin{theorem}[Path locality of exact peeling]\label{thm:dpcm-locality}
Let $R$ be any surviving $r$-clique after round $t$.
If a path $\TPath \notin \Gamma_t(B_t)$, then
\[
\phi_t(R,\TPath)=\phi_{t+1}(R,\TPath).
\]
Consequently,
\[
\deg^{G_{t+1}}_{(s)}(R)
=
\deg^{G_t}_{(s)}(R)
\;-\!\!\!\!\sum_{\TPath \in \Gamma_t(B_t)}\!\!\!\!
\Bigl(\phi_t(R,\TPath)-\phi_{t+1}(R,\TPath)\Bigr).
\]
\end{theorem}

\begin{proof}
If $\TPath \notin \Gamma_t(B_t)$, then neither the path nor the set of active $s$-cliques encoded by it changes during round $t$.
%
Hence its contribution to every surviving $r$-clique remains unchanged.
%
Summing only over affected paths gives the claimed differential form.
\end{proof}

Theorem~\ref{thm:dpcm-locality} is the main algorithmic message of this section:
the global support update problem can be reduced to maintaining a small local state on affected paths.
%
The remaining question is what state is sufficient in different regimes of $r$.

\begin{theorem}[Regime-specific sufficient states]\label{thm:dpcm-states}
Exact differential maintenance on \cpi admits the following sufficient path states.
\begin{enumerate}[leftmargin=*]
    \item \textbf{$r{=}1$.}
    It suffices to maintain whether a path is alive, the number of remaining pivots on the path, and the fixed number of pivots still needed to complete an $s$-clique.
    Hence every support delta is determined by constant-time binomial lookups.

    \item \textbf{$r{=}2$.}
    It suffices to maintain the current \hold and \pivot sets of the path together with the local removal pattern that classifies the update into the three cases studied later.
    Dominant cases admit closed-form support deltas and analytical single-edge splits, avoiding general residual-graph search.

    \item \textbf{$r{\ge}3$.}
    It suffices to maintain the current \hold/\pivot partition, the conflict family induced by removed $r$-cliques, the residual pivot budget, and optional lazy cached clique information for the paths that survive.
    Dead paths can then be identified and handled without recursive splitting, while the remaining paths are updated exactly.
\end{enumerate}
\end{theorem}

\begin{proof}
Each item follows from the path counting semantics in Lemma~\ref{lem:clique_count} and from the later update rules derived for the corresponding regime.
%
For $r{=}1$, the path contribution depends only on how many pivots remain, so counters are sufficient.
%
For $r{=}2$, the contribution of an edge is determined by whether it is hold--hold, hold--pivot, or pivot--pivot, together with the removed pivot pattern.
%
For $r{\ge}3$, the effect of removals is captured by the conflict sets they induce on pivots and by whether the remaining pivot budget still allows a valid path to survive.
\end{proof}

\begin{corollary}[Exactness of DPCM-based peeling]\label{coro:dpcm-exact}
If every affected path applies the exact local update implied by its maintained state in Theorem~\ref{thm:dpcm-states}, then the resulting peeling process computes the exact $(r,s)$-nucleus core number for every $r$-clique.
\end{corollary}

\begin{proof}
The proof is by induction on peeling rounds.
%
Theorem~\ref{thm:dpcm-support} establishes the exact support decomposition at initialization.
%
Theorem~\ref{thm:dpcm-locality} shows that after each peeling step, only affected paths need to be updated.
%
If these local updates are exact, then all surviving $r$-cliques receive exactly their new supports, so the next peeling choice matches the exact peeling process.
\end{proof}

\stitle{Why DPCM is beyond the conference version.}
The conference version introduced \cpi as a compact \emph{representation} of the $s$-clique space.
%
It showed that this representation can replace explicit $s$-clique materialization.
%
The journal-level extension is to separate representation from maintenance.
%
\cpi tells us \emph{where} support is encoded.
%
DPCM tells us \emph{how} exact support changes during peeling.
%
This separation gives a new algorithmic layer.
%
This layer is not present in a purely counting-based view.
%
We no longer maintain support by reconstructing affected clique neighborhoods.
%
Instead, we update path-local sufficient states and recompute only the related differential contributions.

\stitle{Implications for algorithm design.}
Once support maintenance is written in DPCM form, the design problem becomes clear.
%
For each range of $r$, we need to find the smallest sufficient state $\sigma_t(\TPath)$ that still supports exact local deltas.
%
This explains why the three update mechanisms look different.
%
They are still based on the same principle.
%
For $r{=}1$, the sufficient state collapses to counters.
%
For $r{=}2$, the sufficient state becomes an edge-contribution calculus with closed-form cases and analytical splits.
%
For $r{\ge}3$, the sufficient state is conflict-aware and supports dead-path pruning plus lazy residual maintenance.
%
DPCM turns these three optimized algorithms into three exact instantiations of the same maintenance theory.

\stitle{Instantiation map.}
The rest of the paper specializes this idea.
%
Section~\ref{sec:update_cpi} instantiates DPCM for the three main regimes.
%
For $r{=}1$, the path state collapses to two counters and yields an immutable-tree implementation.
%
For $r{=}2$, DPCM becomes a case-based edge contribution calculus with closed-form deltas and analytical path splits.
%
For $r{\ge}3$, DPCM becomes a conflict-aware lazy maintenance scheme.
%
Expensive recursive splitting is used only when the maintained path state shows that a path is still alive.
